% Insert after the one-dimensional theory and before stochastic simulation.
% Requires amsmath and amssymb (already loaded in main-AdvancedCourse.tex).
\section{Multidimensional Master Equations}
\label{sec:multidimensional-master}

So far, we have considered populations whose state is fully characterized by a
single non-negative integer, $n$. In this section we generalize the
construction of the master equation to systems in which several quantities
change in time. Examples include interacting species, the susceptible and
infected individuals of an epidemic, and populations living in different
habitat patches. The guiding principle remains the same: \emph{specify all
possible elementary events and their transition rates, and then account for
probability flowing into and out of each configuration.} We shall develop the
probabilistic description independently of the algorithms used to simulate it.

\subsection{The state vector and the elementary events}
\label{subsec:multi-events}

Consider a system described at time $t$ by the random state vector
\begin{equation}
    \mathbf{N}(t)=(N_1(t),N_2(t),\ldots,N_S(t)),
    \qquad \mathbf{N}(t)\in\mathcal S\subseteq\mathbb{Z}_{\geq 0}^{S}.
    \label{eq:multi-state}
\end{equation}
The components might count individuals of different species, individuals in
various epidemiological compartments, or populations in distinct habitat
patches. The set $\mathcal S$ contains the physically admissible
configurations. The probability that the system has configuration
$\mathbf n=(n_1,\ldots,n_S)$ at time $t$ is written
\begin{equation}
    P(\mathbf n,t)=\Pr\{\mathbf N(t)=\mathbf n\}.
\end{equation}
Normalization requires $\sum_{\mathbf n\in\mathcal S}P(\mathbf n,t)=1$.

Suppose there are $M$ kinds of elementary events, indexed by $k$. Event $k$
changes the state by a fixed \emph{stoichiometric} or \emph{jump} vector
\begin{equation}
    \boldsymbol\nu_k=(\nu_{k1},\ldots,\nu_{kS}),
    \qquad \mathbf n\longrightarrow\mathbf n+\boldsymbol\nu_k,
\end{equation}
and occurs with transition rate (or propensity) $a_k(\mathbf n)\geq 0$.
More precisely, conditional on the current configuration $\mathbf n$,
\begin{equation}
 \Pr\{\text{event $k$ occurs during }[t,t+dt)\mid\mathbf N(t)=\mathbf n\}
 =a_k(\mathbf n)\,dt+o(dt).
 \label{eq:multi-small-time}
\end{equation}
The probability of two or more events during such an interval is $o(dt)$.
We assume here that the rates are time-independent and that the process is
non-explosive. If an event would produce a state outside $\mathcal S$, its
rate must vanish. In this way, boundary conditions are encoded in the event
rates themselves.

The collection $\{\boldsymbol\nu_k,a_k(\mathbf n)\}_{k=1}^M$ defines a
continuous-time Markov jump process: once its present configuration is known,
its future evolution does not depend on how it arrived there. Notice that a
single elementary event may change \emph{more than one} component of the
state vector at once. This is a crucial difference from a one-dimensional
birth--death process, even though we still describe dynamics as a sequence of
individual jumps.

\subsection{Derivation of the general master equation}
\label{subsec:multi-derivation}

Let $A(\mathbf n)=\sum_{k=1}^M a_k(\mathbf n)$ denote the total rate of
leaving configuration $\mathbf n$. During a sufficiently short interval
$dt$, a system found at $\mathbf n$ either remains there or is reached from
one of the configurations $\mathbf n-\boldsymbol\nu_k$ through event $k$.
Thus, up to terms smaller than $dt$,
\begin{align}
 P(\mathbf n,t+dt)={}&
   [1-A(\mathbf n)dt]P(\mathbf n,t) \nonumber\\
 &+dt\sum_{k=1}^{M}
       a_k(\mathbf n-\boldsymbol\nu_k)
       P(\mathbf n-\boldsymbol\nu_k,t)+o(dt).
 \label{eq:multi-bookkeeping}
\end{align}
The first term accounts for systems which remain in $\mathbf n$; each term
in the sum is an incoming probability flux. Subtracting $P(\mathbf n,t)$,
dividing by $dt$, and taking $dt\to0$ gives the \emph{multidimensional
master equation}:
\begin{equation}
 \boxed{\displaystyle
 \frac{\partial P(\mathbf n,t)}{\partial t}
 =\sum_{k=1}^{M}\left[
    a_k(\mathbf n-\boldsymbol\nu_k)
    P(\mathbf n-\boldsymbol\nu_k,t)
    -a_k(\mathbf n)P(\mathbf n,t)
  \right].}
 \label{eq:multi-master}
\end{equation}
Whenever $\mathbf n-\boldsymbol\nu_k$ is not an admissible state, the
corresponding incoming term is understood to be zero. As in the
one-dimensional case, this is not a single ordinary differential equation:
it is a coupled system of equations, one for every configuration in
$\mathcal S$.

\paragraph{Recovering the one-dimensional birth--death equation.}
For $S=1$, choose two events with
$\nu_1=+1$, $a_1(n)=g_n$ and $\nu_2=-1$,
$a_2(n)=r_n$. Equation~\eqref{eq:multi-master} becomes
\begin{equation}
 \frac{dP(n,t)}{dt}
 =g_{n-1}P(n-1,t)+r_{n+1}P(n+1,t)
  -(g_n+r_n)P(n,t),
 \label{eq:multi-recovery}
\end{equation}
precisely the birth--death master equation derived earlier (with the
appropriate treatment at $n=0$). The higher-dimensional formulation is
therefore a direct generalization, not a different kind of stochastic model.

\subsection{A worked example: stochastic predator--prey interactions}
\label{subsec:multi-predator-prey}

Consider two populations: prey, whose number is $n$, and predators, whose
number is $p$. We begin with a deliberately simple model. Prey reproduce
independently at per-capita rate $b$; predator--prey encounters occur at
rate $a$ per potential pair and each successful encounter consumes one prey
and produces one additional predator; predators die at per-capita rate $d$.
The state is $\mathbf n=(n,p)$, with $n,p=0,1,\ldots$.

\begin{center}
\begin{tabular}{llll}
\hline
Elementary event & State change & Jump vector & Rate \\
\hline
Prey reproduction & $(n,p)\to(n+1,p)$ & $(+1,0)$ & $bn$ \\
Predation and predator reproduction & $(n,p)\to(n-1,p+1)$ & $(-1,+1)$ & $anp$ \\
Predator death & $(n,p)\to(n,p-1)$ & $(0,-1)$ & $dp$ \\
\hline
\end{tabular}
\end{center}

The term $anp$ is a mass-action encounter rate at fixed habitat area. The
second event is particularly instructive: the disappearance of one prey and
the appearance of one predator are \emph{one single elementary event}, not
two independent events. All three rates are non-negative and automatically
vanish when a required individual is absent.

We now write the equation for $P(n,p,t)$ by accounting separately for the
three ways in which the configuration $(n,p)$ can be reached:
\begin{itemize}
 \item A prey birth from $(n-1,p)$, at rate $b(n-1)$.
 \item A predation event from $(n+1,p-1)$, at rate $a(n+1)(p-1)$.
 \item A predator death from $(n,p+1)$, at rate $d(p+1)$.
\end{itemize}
The probability also decreases whenever one of the three events occurs in
$(n,p)$. The master equation is therefore
\begin{align}
 \frac{\partial P(n,p,t)}{\partial t}={}&
       b(n-1)P(n-1,p,t) \nonumber\\
 &+a(n+1)(p-1)P(n+1,p-1,t) \nonumber\\
 &+d(p+1)P(n,p+1,t) \nonumber\\
 &-[bn+anp+dp]P(n,p,t).
 \label{eq:multi-predator-master}
\end{align}
Terms with negative arguments are zero. At $n=0$ or $p=0$, some terms
vanish: there is no predation without prey or predators, and neither
population can become negative. If either species becomes extinct in this
minimal model, it cannot spontaneously reappear. More realistic versions
could include immigration, density dependence, or other processes; these
would simply add further rows to the event table and terms to the master
equation.

\subsection{Equations for averages and the deterministic description}
\label{subsec:multi-moments}

The master equation describes the entire probability distribution. Often,
however, we are also interested in average abundances. For any function
$f(\mathbf n)$, multiply Eq.~\eqref{eq:multi-master} by $f(\mathbf n)$,
sum over all configurations, and shift indices in the incoming terms. This
gives
\begin{equation}
 \boxed{\displaystyle
 \frac{d}{dt}\langle f(\mathbf N)\rangle
 =\left\langle\sum_{k=1}^{M}a_k(\mathbf N)
    [f(\mathbf N+\boldsymbol\nu_k)-f(\mathbf N)]\right\rangle.}
 \label{eq:multi-observable}
\end{equation}
Here $\langle\cdot\rangle$ denotes an ensemble average. In particular,
with $f(\mathbf n)=n_i$ we obtain the exact moment equation
\begin{equation}
 \boxed{\displaystyle
 \frac{d\langle N_i\rangle}{dt}
   =\sum_{k=1}^{M}\nu_{ki}\langle a_k(\mathbf N)\rangle.}
 \label{eq:multi-first-moment}
\end{equation}
This result provides a direct bridge from elementary stochastic events to
macroscopic rates of change.

For the predator--prey example, it gives
\begin{align}
 \frac{d\langle n\rangle}{dt}
      &=b\langle n\rangle-a\langle np\rangle,\\
 \frac{d\langle p\rangle}{dt}
      &=a\langle np\rangle-d\langle p\rangle.
 \label{eq:multi-predator-means}
\end{align}
These equations are \emph{exact}, but they are not closed: the average
$\langle np\rangle$ depends on the joint distribution of prey and predator
abundances. Indeed,
\begin{equation}
 \langle np\rangle=\langle n\rangle\langle p\rangle
                       +\operatorname{Cov}(n,p).
 \label{eq:multi-covariance}
\end{equation}
Replacing $\langle np\rangle$ by
$\langle n\rangle\langle p\rangle$ neglects this covariance and produces
the familiar deterministic mass-action equations for mean abundances. This
is a \emph{moment-closure approximation}, not in general an exact identity.
Stochastic population dynamics can generate correlations even when the
initial abundances are independent.

A deterministic limit can also be motivated by introducing a system-size
parameter $\Omega$ (for example, habitat area) and density variables
$x=n/\Omega$, $y=p/\Omega$. For a fixed encounter coefficient $\alpha$,
we then write the predation rate as $\alpha np/\Omega$ so that all
macroscopic rates scale with $\Omega$. If relative fluctuations become
negligible as $\Omega$ increases, the limiting densities satisfy
\begin{equation}
 \frac{dx}{dt}=bx-\alpha xy,\qquad
 \frac{dy}{dt}=\alpha xy-dy.
 \label{eq:multi-deterministic-limit}
\end{equation}
The conditions under which such an approximation is accurate deserve
careful attention. In particular, extinction events and fluctuations near
small populations are precisely among the phenomena that may be lost in a
deterministic approximation. A systematic treatment of the
large-system-size approximation can be developed separately (see Appendix \ref{App:LargeN-expansion}).

\subsection{Why the multidimensional formulation matters}
\label{subsec:multi-why}

The same construction works whether components of the state vector refer to
species, infection classes, or locations. In a spatially explicit
metapopulation model, for example, $N_i$ might be the abundance in habitat
patch $i$. A dispersal event from patch $i$ to patch $j$ then has jump
vector $-\mathbf e_i+\mathbf e_j$, where $\mathbf e_i$ is the unit vector
for patch $i$. Such an event simultaneously decreases one local population
and increases another. All spatial information is encoded in the allowed
events and their rates, not in a change of the master-equation principle.

The essential distinction is thus between \emph{specifying the stochastic
model} and \emph{choosing how to study it}. The master equation provides a
complete probabilistic description. Stationary distributions, moment
equations, approximations, and stochastic simulations are different ways of
extracting information from that same model.

\subsection{Exercises}
\label{subsec:multi-exercises}
\begin{enumerate}
 \item \textbf{Independent populations.} Suppose two species reproduce and die independently, with transition rates $b_i n_i$ and $d_i n_i$ for species $i=1,2$. Write the four jump vectors and derive the two-dimensional master equation. Under what initial conditions does its joint distribution factorize at later times?
 \item \textbf{Predation and boundaries.} Starting from Eq.~\eqref{eq:multi-predator-master}, write explicitly the equations for $P(0,p,t)$ and $P(n,0,t)$. Explain which events remain possible at these boundaries.
 \item \textbf{Moment closure.} Derive an equation for $d\langle np\rangle/dt$ using Eq.~\eqref{eq:multi-observable}. Identify the higher-order moments that appear. Why does an equation for the average not generally close on its own?
 \item \textbf{Dispersal between two sites.} Let $(n_1,n_2)$ describe individuals in two habitat patches. Individuals disperse from patch 1 to patch 2 at rate $m_{12}n_1$, and in the reverse direction at rate $m_{21}n_2$. Derive the master equation and show that $n_1+n_2$ is conserved along every trajectory.
\end{enumerate}
